% ADMD Lecture 14. Compile: latexmk -xelatex lecture14.tex
\documentclass[10pt,aspectratio=169,t]{beamer}
\usetheme{upbminimal}

\course{Application Development for Mobile Devices}
\decklabel{Lecture 14}
\title{UI polish and testing}
\subtitle{Theming, layout, forms, motion, accessibility, and tests}
\author{Dinu-Ștefan Rusu}
\institute{FILS, Universitatea Politehnica București}
\date{Autumn 2026}

\begin{document}

\titleframe

\objectivesframe{
  \cell[\faPalette]{One theme, two modes}{Build the app's colors and type once with ThemeData, and get dark mode for free.}
  \cell[\faKeyboard]{Forms that behave}{Validators, focus order, keyboard types, and error messages a user can act on.}
  \cell[\faUniversalAccess]{Reachable by everyone}{Semantics, contrast, tap targets, text scaling, and localized strings.}
  \cell[\faVial]{Tests you trust}{Unit, widget, bloc, golden and integration tests, all running on every pull request.}
}

\divider{Part 1 · Theming}{Making it look like one app}[ThemeData, Material 3, dark mode, and responsive layout]

\begin{frame}[fragile,eyebrow=Theming]{Colors do not belong inside widgets}
  \begin{columns}[T]
    \begin{column}{0.55\textwidth}
\begin{dart}
// Repeated in forty widgets:
Container(color: const Color(0xFF3F51B5))
Text('Total', style: TextStyle(
  fontSize: 16, color: Colors.black87))

// Declared once, read everywhere:
final cs = Theme.of(context).colorScheme;
Container(color: cs.primaryContainer)
Text('Total', style: tt.titleMedium)
\end{dart}
    \end{column}
    \begin{column}{0.41\textwidth}
      \lead{The first version has no dark mode.}{Every hard-coded color is a value that cannot change with the platform, the user's settings, or next year's brand.}
      \muted{It is also unsearchable. Nobody can answer "which blue is the button" without reading forty files.}
    \end{column}
  \end{columns}
  \begin{takeaway}{Read a role, never write a color.}
    A widget that names \code{cs.primary} keeps working when the palette, the mode or the brand changes.
  \end{takeaway}
  \note{Ask what happens to the top version when the phone switches to dark mode. Nothing happens: that is the point. Say that theming is not decoration, it is removing duplicated constants.}
\end{frame}

\begin{frame}[fragile,eyebrow=Theming]{ThemeData and ColorScheme.fromSeed}
\begin{dart}
ColorScheme _scheme(Brightness b) => ColorScheme.fromSeed(
      seedColor: const Color(0xFF3F51B5), brightness: b);

ThemeData appTheme(Brightness b) => ThemeData(colorScheme: _scheme(b));

// Read it in any widget, never re-declare it:
final cs = Theme.of(context).colorScheme;
final tt = Theme.of(context).textTheme;
\end{dart}
  \begin{takeaway}{One seed color produces the whole palette.}
    Material 3 derives every role from it and keeps the contrast between a color and its \code{on} pair correct in both modes.
  \end{takeaway}
  \note{Point at the seed: it is the only color the designer picks. Show that the generated palette is not the seed repeated, it is a set of tones. Say the theme lives in one file, for example lib/theme/app\_theme.dart. Mention that Material has moved into the separate material\_ui package, that the in-SDK library is deprecated from the November 2026 release, and that dart fix performs the migration.}
\end{frame}

\begin{frame}[eyebrow=Theming]{The color roles you will actually use}
  \vskip 1mm
  \begin{tabular}{@{}lp{62mm}l@{}}
    \toprule
    \thead{Role} & \thead{What it paints} & \thead{Text on it} \\
    \midrule
    primary             & the main action: a filled button    & onPrimary \\
    primaryContainer    & a tinted area for that action       & onPrimaryContainer \\
    surface             & page background, cards, sheets      & onSurface \\
    secondary, tertiary & accents that are not the main action & onSecondary \\
    error               & validation and destructive actions  & onError \\
    outline             & borders, dividers, disabled edges   & \\
    \bottomrule
  \end{tabular}
  \vskip 2mm
  \muted{Every role has an \code{on} pair. Paint with one, write with the other, and contrast is already handled.}
  \note{Stress the on-pair rule, because it is the single habit that prevents unreadable text in dark mode. Ask the room which role a delete button uses. The answer is error, not a red hex code.}
\end{frame}

\begin{frame}[fragile,eyebrow=Theming]{Dark mode is a second ThemeData}
\begin{dart}
MaterialApp(
  theme: appTheme(Brightness.light),
  darkTheme: appTheme(Brightness.dark),
  themeMode: ThemeMode.system,   // or the user's saved choice
  home: const HomeScreen(),
);

// Rare, and usually a smell:
final isDark = Theme.of(context).brightness == Brightness.dark;
\end{dart}
  \begin{takeaway}{You do not design dark mode twice.}
    If widgets read roles instead of colors, the same widget tree is already correct in both modes.
  \end{takeaway}
  \note{Say that themeMode should be persisted with shared\_preferences from Lecture 9 once the user picks one. Warn that images and illustrations with baked-in white backgrounds are the one thing dark mode cannot fix.}
\end{frame}

\begin{frame}[fragile,eyebrow=Theming]{Typography and component themes}
\begin{dart}
ThemeData(
  colorScheme: _scheme(b),
  fontFamily: 'Inter',
  appBarTheme: const AppBarTheme(centerTitle: true),
  inputDecorationTheme: const InputDecorationTheme(
    border: OutlineInputBorder(),
  ),
  filledButtonTheme: FilledButtonThemeData(
    style: FilledButton.styleFrom(minimumSize: const Size.fromHeight(48)),
  ),
);
\end{dart}
  \vskip 1mm
  \muted{The scale runs from \code{displayLarge} to \code{labelSmall}. Most screens need only three of its roles.}
  \note{Component themes are how you stop copying the same button style. Set the minimum button height here once and the tap-target rule from Part 4 is satisfied everywhere. Mention that a custom font is declared in pubspec.yaml under flutter: fonts.}
\end{frame}

\begin{frame}[fragile,eyebrow=Layout]{MediaQuery: what the window is doing}
  \begin{columns}[T]
    \begin{column}{0.55\textwidth}
\begin{dart}
final size = MediaQuery.sizeOf(context);
final pad = MediaQuery.paddingOf(context);
final keyboard =
    MediaQuery.viewInsetsOf(context).bottom;
final scaler =
    MediaQuery.textScalerOf(context);
\end{dart}
    \end{column}
    \begin{column}{0.41\textwidth}
      \lead{Ask for one metric, not all of them.}{\code{MediaQuery.of(context)} rebuilds your widget when any metric changes, including the keyboard opening. The \code{sizeOf} family rebuilds only on the value you read.}
      \muted{\code{viewInsets} is the keyboard. \code{padding} is the notch and the home indicator.}
    \end{column}
  \end{columns}
  \begin{takeaway}{MediaQuery describes the window, not your widget.}
    It reports what the system is doing to the app: the size, the notch, the keyboard, and the user's text scale.
  \end{takeaway}
  \note{Show what happens when a widget reads MediaQuery.of and the keyboard slides up: the whole subtree rebuilds on every animation frame. This is the most common accidental jank in a form screen.}
\end{frame}

\begin{frame}[fragile,eyebrow=Layout]{LayoutBuilder, breakpoints, and SafeArea}
\begin{dart}
LayoutBuilder(builder: (context, constraints) {
  if (constraints.maxWidth < 600) return const TodoList();
  return const Row(children: [
    Expanded(flex: 2, child: TodoList()),
    Expanded(flex: 3, child: TodoDetail()),
  ]);
});

SafeArea(child: body);  // keeps content clear of notch and gestures
\end{dart}
  \begin{takeaway}{Break on the width you were given, not on the device.}
    A phone in landscape, a tablet, and a split-screen window can all be the same width. Two breakpoints, near 600 and near 840, cover almost everything.
  \end{takeaway}
  \note{MediaQuery describes the window, LayoutBuilder describes the box this widget was actually handed. Inside a side panel they disagree, and LayoutBuilder is the one that is right. Show the app in a resized desktop window if the emulator allows it.}
\end{frame}

\divider{Part 2 · Forms}{Getting data out of a person}[Form, TextFormField, validators, focus, and error display]

\begin{frame}[fragile,eyebrow=Forms]{The shape of a Flutter form}
\begin{dart}
final _formKey = GlobalKey<FormState>();

Form(
  key: _formKey,
  autovalidateMode: AutovalidateMode.onUserInteraction,
  child: Column(children: [emailField, passwordField, submitButton]),
);

if (_formKey.currentState!.validate()) await _signIn();
\end{dart}
  \begin{takeaway}{The Form is a coordinator, not a container.}
    It finds every descendant field, runs their validators, and reports one boolean. The fields keep their own state.
  \end{takeaway}
  \note{Explain that GlobalKey is the exception to the rule that you do not reach into other widgets. Point out autovalidateMode: onUserInteraction, which avoids showing errors on an empty form the user has not touched yet.}
\end{frame}

\begin{frame}[fragile,eyebrow=Forms]{Validators that say what to fix}
  \begin{columns}[T]
    \begin{column}{0.55\textwidth}
\begin{dart}
String? validateEmail(String? v) {
  final s = v?.trim() ?? '';
  if (s.isEmpty) return 'Enter your email';
  if (!s.contains('@')) {
    return 'This address is missing an @';
  }
  return null;   // null means valid
}
\end{dart}
    \end{column}
    \begin{column}{0.41\textwidth}
      \lead{Return null when the value is fine.}{Any other string becomes the red line under the field. One message, naming the one thing to change.}
      \muted{"Invalid input" tells the user nothing. Trim before you check, and never reject a password for containing a space.}
    \end{column}
  \end{columns}
  \begin{takeaway}{A validator is a pure function.}
    A string in, a string or null out. That makes it the easiest thing in the whole app to unit test, as Part 5 shows.
  \end{takeaway}
  \note{Ask the room for the worst validation message they have met. Usually it is a password rule that is not stated up front. Say that a validator is a pure function, which makes it the easiest thing in the whole app to unit test.}
\end{frame}

\begin{frame}[fragile,eyebrow=Forms]{Keyboard, focus order, and submit}
\begin{dart}
TextFormField(
  controller: _email,
  focusNode: _emailFocus,
  keyboardType: TextInputType.emailAddress,
  textInputAction: TextInputAction.next,
  onFieldSubmitted: (_) => _passwordFocus.requestFocus(),
  decoration: const InputDecoration(labelText: 'Email'),
);

// In dispose(): _email.dispose(); _emailFocus.dispose();
\end{dart}
  \begin{takeaway}{Dispose everything you create.}
A controller and a FocusNode belong to your State object, and leak if forgotten.
  \end{takeaway}
  \note{Every controller and every FocusNode you create must be disposed, or the State leak shows up in DevTools memory. Demonstrate the difference between TextInputAction.next and done on a real keyboard: it is a two-line change that people notice immediately.}
\end{frame}

\begin{frame}[eyebrow=Forms]{Errors people can act on}
  \vskip 2mm
  \begin{grid}[3]
    \cell[\faExclamationCircle]{Errors under the field}{The message sits where the mistake is. Nobody hunts for the wrong field.}
    \cell[\faServer]{Server errors}{"Email already registered" is not a validator result. Show it near the button.}
    \cell[\faLock]{Disable submit in flight}{Keep a bool in the bloc state. A double tap must not create two accounts.}
    \cell[\faUndo]{Keep what they typed}{On failure, keep every field. Retyping after a network blink loses users.}
    \cell[\faArrowUp]{Scroll to the error}{On a long form the field can be off screen. \code{Scrollable.ensureVisible} helps.}
    \cell[\faCheckCircle]{Say what succeeded}{A form that submits silently looks broken. One SnackBar is enough.}
  \end{grid}
  \note{These six rules cost about twenty lines and separate a student project from a shipped app. Ask which one they have personally been annoyed by this week.}
\end{frame}

\divider{Part 3 · Motion}{Animation you get for free}[Implicit animations, Hero, and knowing when to stop]

\begin{frame}[eyebrow=Motion]{The implicit animation family}
  \vskip 2mm
  \begin{grid}[3]
    \cell[\faExpand]{AnimatedContainer}{Change size, color, padding or decoration and it tweens between old and new.}
    \cell[\faEye]{AnimatedOpacity}{Fade in or out. An opacity of zero still lays out and still hits tests.}
    \cell[\faRandom]{AnimatedSwitcher}{Cross-fades between two children when the child's key changes.}
    \cell[\faAlignLeft]{AnimatedAlign}{Move a child inside its parent. So does AnimatedPadding.}
    \cell[\faFont]{AnimatedDefaultTextStyle}{Grow, shrink or recolor text as a state changes.}
    \cell[\faMagic]{TweenAnimationBuilder}{Any value you can tween, animated once, with no State of your own.}
  \end{grid}
  \note{The rule of thumb: if the animation runs because a value changed, it is implicit. If it runs because time passed, it is explicit. Most product animation is the first kind, which is why you rarely need an AnimationController. Say that explicit animations, driven by a controller and a ticker, are for motion you must repeat or reverse.}
\end{frame}

\begin{frame}[fragile,eyebrow=Motion]{Two widgets, one state change}
\begin{dart}
const d = Duration(milliseconds: 200);   // one place, whole app
AnimatedContainer(
  duration: d, curve: Curves.easeOut,
  color: done ? cs.primaryContainer : cs.surface,
  padding: EdgeInsets.all(done ? 8 : 16),
  child: const TodoTitle(),
);
AnimatedSwitcher(duration: d,
  child: loading ? const Spinner(key: ValueKey('spin'))
                 : TodoList(key: ValueKey(items.length)));
\end{dart}
  \begin{takeaway}{The key is what makes AnimatedSwitcher work.}
    Without a different key, Flutter reuses the same element and there is nothing to cross-fade.
  \end{takeaway}
  \note{Remove the ValueKey from the AnimatedSwitcher children live and show that nothing animates: Flutter sees the same widget type and reuses the element. That single detail is the most common bug report about AnimatedSwitcher.}
\end{frame}

\begin{frame}[fragile,eyebrow=Motion]{Hero: one element across two routes}
  \begin{columns}[T]
    \begin{column}{0.55\textwidth}
\begin{dart}
// List screen
Hero(
  tag: 'todo-${todo.id}',
  child: TodoThumbnail(todo),
);

// Detail screen, same tag
Hero(tag: 'todo-${todo.id}', child: big);
\end{dart}
    \end{column}
    \begin{column}{0.41\textwidth}
      \lead{The tag is the contract.}{Flutter finds the widget carrying the same tag on both routes and flies it between the two positions.}
      \muted{A tag must be unique on each screen. Use the record's id, not the list index.}
    \end{column}
  \end{columns}
  \begin{takeaway}{Hero borrows the route transition.}
    No controller, no duration. The only thing you own is the tag, and duplicate tags throw.
  \end{takeaway}
  \note{Hero needs no controller and no duration: it borrows the route transition. Show it with a thumbnail that becomes a header image. Warn that text inside a Hero can flash unstyled unless it sits under a Material widget.}
\end{frame}

\begin{frame}[fragile,eyebrow=Motion]{When to stop animating}
\begin{dart}
if (MediaQuery.disableAnimationsOf(context)) return child;
\end{dart}
  \vskip 1mm
  \begin{tabular}{@{}lp{48mm}p{58mm}@{}}
    \toprule
    \thead{Length} & \thead{Use it for} & \thead{Sign you got it wrong} \\
    \midrule
    100 to 150 ms & a ripple, a checkbox, a fade & the change is invisible \\
    200 to 300 ms & cards, sheets, list items    & the app feels slow to a fast user \\
    300 to 400 ms & a full route transition      & the user waits for the animation \\
    \bottomrule
  \end{tabular}
  \begin{takeaway}{Motion explains a change, it does not announce it.}
If the user notices the animation instead of the result, it is too long. Never block input until it ends.
  \end{takeaway}
  \note{Reduce Motion is a real accessibility setting on both platforms and it reaches you through MediaQuery. For some users animation causes nausea, so this is not a preference to ignore. Say that motion should explain a change, not announce the developer.}
\end{frame}

\divider{Part 4 · Reach}{Accessibility and localization}[Semantics, contrast, tap targets, text scaling, ARB files]

\begin{frame}[eyebrow=Access]{Six things that lock people out}
  \vskip 2mm
  \begin{grid}[3]
    \cell[\faVolumeUp]{Unlabeled controls}{An icon button with no label is announced as "button" and nothing else by TalkBack and VoiceOver.}
    \cell[\faLowVision]{Low contrast}{Gray on gray fails outdoors, on a cheap panel, and for anyone over forty.}
    \cell[\faHandPointer]{Tiny targets}{A 24 dp icon is a 24 dp target unless you say otherwise. Fingers are not cursors.}
    \cell[\faTextHeight]{Fixed heights}{A row with a hard-coded height overflows the moment the user enlarges text.}
    \cell[\faPalette]{Color as the only signal}{Red text alone does not say "error" to a person who cannot see red. Add an icon or a word.}
    \cell[\faKeyboard]{No focus order}{On a keyboard or a switch device, tab order is your layout order. Check it once.}
  \end{grid}
  \note{Turn TalkBack on for thirty seconds in the lecture and navigate one screen. It is the fastest way to make this section concrete. Say that accessibility work is also plain quality work: it fixes bugs sighted users hit too.}
\end{frame}

\begin{frame}[fragile,eyebrow=Access]{Semantics: what the screen reader says}
  \begin{columns}[T]
    \begin{column}{0.55\textwidth}
\begin{dart}
IconButton(
  icon: const Icon(Icons.delete),
  tooltip: 'Delete todo',   // also a label
  onPressed: _delete,
);

Semantics(
  label: '3 of 7 done',
  child: ExcludeSemantics(child: chart),
);
\end{dart}
    \end{column}
    \begin{column}{0.41\textwidth}
      \lead{Describe the meaning, not the pixels.}{A chart is a picture to the framework. \code{Semantics} replaces it with the sentence a person needs.}
      \muted{\code{MergeSemantics} joins a row into one announcement. \code{liveRegion: true} makes a status line speak when it changes.}
    \end{column}
  \end{columns}
  \begin{takeaway}{The framework guesses, and you confirm.}
    Flutter builds a semantics tree on its own. Your job is every place where the pixels do not carry the meaning.
  \end{takeaway}
  \note{Point out that tooltip on an IconButton gives you a semantic label for free, so there is no excuse for an unlabeled icon. Open the Semantics debugger in DevTools to show the tree the reader actually walks.}
\end{frame}

\begin{frame}[eyebrow=Access]{Contrast, targets, and text that grows}
  \vskip 1mm
  \begin{tabular}{@{}lp{30mm}p{60mm}@{}}
    \toprule
    \thead{Rule} & \thead{Target} & \thead{How to satisfy it} \\
    \midrule
    Body text contrast   & 4.5 to 1        & an on role over its own surface role \\
    Large text and icons & 3 to 1          & the same, plus a design-tool check \\
    Tap target           & 48 dp, 44 pt    & a minimum button size in the theme \\
    Text scaling         & up to 2 times   & no fixed heights, let rows wrap \\
    Focus order          & follows reading & widget order equals visual order \\
    \bottomrule
  \end{tabular}
  \begin{takeaway}{Test with the phone's own settings.}
    Set the font size to the largest step, turn a screen reader on, and open every screen. Two minutes, and it finds most of these.
  \end{takeaway}
  \note{The two contrast numbers come from WCAG level AA and every serious design review uses them. Note that Material's own components already meet the 48 dp target, so the failures are almost always custom widgets.}
\end{frame}

\begin{frame}[fragile,eyebrow=Locale]{Strings live in ARB files, not in widgets}
\begin{yaml}
# pubspec.yaml
flutter:
  generate: true

# l10n.yaml
arb-dir: lib/l10n
template-arb-file: app_en.arb
output-localization-file: app_localizations.dart
\end{yaml}
  \vskip 2mm
  \muted{Add \code{flutter\_localizations} and \code{intl} as dependencies. Building the app regenerates the \code{AppLocalizations} class from the ARB files. Pass \code{AppLocalizations.localizationsDelegates} and \code{supportedLocales} to \code{MaterialApp}, then read a string with \code{AppLocalizations.of(context)}.}
  \note{ARB is a JSON file of key and message pairs, one per locale. The generated class means a missing or misspelled key is a compile error, not a blank label at runtime. Say that translators are given the ARB file, never the source.}
\end{frame}

\begin{frame}[fragile,eyebrow=Locale]{Plurals, dates, and right to left}
\begin{codeblock}
// lib/l10n/app_en.arb
"todoCount": "{count, plural, =0{No todos} =1{1 todo} other{{count} todos}}",
"@todoCount": { "placeholders": { "count": { "type": "int" } } }

// In the widget
Text(AppLocalizations.of(context)!.todoCount(items.length));
Text(DateFormat.yMMMd(locale).format(todo.due));

// Right to left: start and end, never left and right
padding: const EdgeInsetsDirectional.only(start: 16),
\end{codeblock}
  \vskip 1mm
  \muted{Some languages have more than two plural forms, which is why you never build a sentence by concatenation.}
  \note{Show what "1 todos" looks like and why string concatenation cannot be translated: word order changes between languages. For right to left, wrap a screen in Directionality with rtl and watch which paddings stay on the wrong side.}
\end{frame}

\divider{Part 5 · Testing}{Proving it still works}[Unit, widget, bloc, golden and integration tests]

\begin{frame}[eyebrow=Testing]{The pyramid, and what each layer buys}
  \vskip 1mm
  {\usebeamerfont{small}\begin{tabular}{@{}llp{20mm}p{52mm}@{}}
    \toprule
    \thead{Layer} & \thead{Runs on} & \thead{Speed} & \thead{What it proves} \\
    \midrule
    Unit        & the Dart VM, no UI   & milliseconds & a validator, a model, a repository \\
    Widget      & a headless tree      & milliseconds & a screen builds and reacts to a tap \\
    Bloc        & the Dart VM          & milliseconds & an event produces these states \\
    Golden      & a headless tree      & fast         & the pixels did not change by accident \\
    Integration & a device or emulator & seconds      & the whole app, plugins and network \\
    \bottomrule
  \end{tabular}}
  \vskip 2mm
  \begin{takeaway}{Write many cheap tests and a few expensive ones.}
    Every layer catches a bug the layer below cannot see, and costs more to write and to run.
  \end{takeaway}
  \note{Ask how many have written a test for a UI before. Usually nobody. Say the goal is not coverage as a number, it is that a refactor tomorrow tells you within a minute whether it broke something.}
\end{frame}

\begin{frame}[fragile,eyebrow=Testing]{Unit tests: the cheapest ones first}
\begin{dart}
// test/validators_test.dart
void main() {
  group('validateEmail', () {
    test('rejects an empty value', () {
      expect(validateEmail(''), 'Enter your email');
    });
    test('accepts a normal address', () {
      expect(validateEmail('ana@example.com'), isNull);
    });
  });
}
\end{dart}
  \vskip 2mm
  \muted{Files end in \code{\_test.dart} and live under \code{test/}. Run them all with \code{flutter test}.}
  \note{Note the naming: the test name is a sentence about behavior, not about the method. Point out that pure functions like validators, copyWith and fromJson are where the ratio of value to effort is highest.}
\end{frame}

\begin{frame}[fragile,eyebrow=Testing]{Widget tests: testWidgets, find, pump}
\begin{dart}
testWidgets('adding a todo shows it in the list', (tester) async {
  await tester.pumpWidget(const MaterialApp(home: TodoScreen()));

  await tester.enterText(find.byType(TextField), 'Buy milk');
  await tester.tap(find.byIcon(Icons.add));
  await tester.pump();                       // build one new frame

  expect(find.text('Buy milk'), findsOneWidget);
  expect(find.byType(TodoTile), findsNWidgets(1));
});
\end{dart}
  \vskip 1mm
  \muted{Finders: \code{find.text}, \code{find.byType}, \code{find.byIcon}, \code{find.byKey}. Matchers: \code{findsOneWidget}, \code{findsNothing}, \code{findsNWidgets(n)}.}
  \note{There is no device and no screen here: the widget tree is built in memory, which is why these run in milliseconds. Say that find.byKey is the stable choice, because visible text changes when the app is localized.}
\end{frame}

\begin{frame}[fragile,eyebrow=Testing]{pump, pumpAndSettle, and the trap}
  \begin{columns}[T]
    \begin{column}{0.55\textwidth}
\begin{dart}
await tester.pump();
// one frame

await tester.pump(
  const Duration(milliseconds: 300));
// advance time by 300 ms

await tester.pumpAndSettle();
// frames until no animation is left
\end{dart}
    \end{column}
    \begin{column}{0.41\textwidth}
      \lead{Test time does not pass on its own.}{Nothing animates, no timer fires and no Future completes until you pump. That is a feature: the test is deterministic.}
      \begin{warning}
        \code{pumpAndSettle} never returns if something animates forever, such as a spinner. Pump a fixed duration instead.
      \end{warning}
    \end{column}
  \end{columns}
  \note{The hanging test with a CircularProgressIndicator on screen is the single most common failure students hit. Show it once, then show the fix. Also mention tester.pumpWidget wraps a MaterialApp because most widgets need Directionality and a Theme.}
\end{frame}

\begin{frame}[fragile,eyebrow=Testing]{Faking the outside world with mocktail}
\begin{dart}
class MockTodoRepository extends Mock implements TodoRepository {}

late MockTodoRepository repo;

setUp(() {
  repo = MockTodoRepository();
  when(() => repo.load()).thenAnswer((_) async => [Todo(id: 1)]);
});

test('the cubit asks the repository once', () async {
  await TodoCubit(repo).load();
  verify(() => repo.load()).called(1);
});
\end{dart}
  \note{A test that calls the real network is not a test, it is a flaky report about someone else's server. Explain that mocktail needs no code generation, which is why it is preferred here over the older alternative. Mention registerFallbackValue when an argument matcher is used on a custom type.}
\end{frame}

\begin{frame}[fragile,eyebrow=Testing]{Testing a bloc: back to Lecture 6}
\begin{dart}
blocTest<TodoBloc, TodoState>(
  'emits loading then loaded when TodosRequested is added',
  build: () => TodoBloc(repo),          // repo is a mock, see previous slide
  act: (bloc) => bloc.add(const TodosRequested()),
  expect: () => [
    const TodoState.loading(),
    TodoState.loaded([Todo(id: 1)]),
  ],
);
\end{dart}
  \begin{takeaway}{This is why state classes carry value equality.}
    \code{expect} compares the emitted states to your list. Without Equatable, two identical states are not equal and the test fails.
  \end{takeaway}
  \note{Tie this back to the state design rules from Lecture 6. A bloc test is a description of a feature in three lines: given this, when that event arrives, these states come out. It also documents the feature better than a comment.}
\end{frame}

\begin{frame}[fragile,eyebrow=Testing]{Golden tests: the pixels themselves}
\begin{dart}
testWidgets('TodoTile matches its golden', (tester) async {
  await tester.pumpWidget(wrap(const TodoTile(title: 'Buy milk')));
  await expectLater(
    find.byType(TodoTile),
    matchesGoldenFile('goldens/todo_tile.png'),
  );
});
\end{dart}
  \vskip 2mm
  \muted{Create or update the reference images with \code{flutter test --update-goldens}, then review them in the pull request like any other change.}
  \begin{takeaway}{Commit the golden files.}
    The diff in the review is a picture, which is the only test output a designer can read.
  \end{takeaway}
  \note{Warn that goldens are only stable because the test environment uses a fixed test font, not the system font. Keep goldens for small reusable widgets. A golden of a whole screen fails on every unrelated change and gets ignored.}
\end{frame}

\begin{frame}[fragile,eyebrow=Testing]{integration\_test, and running all of it in CI}
\begin{shell}
flutter analyze
flutter test                        # unit, widget, bloc, golden
flutter test integration_test       # on a booted device or emulator
\end{shell}
  \vskip 2mm
  \muted{An integration test starts with \code{IntegrationTestWidgetsFlutterBinding.ensureInitialized()} and then uses the same \code{testWidgets} API, except that the real app runs on a real device with real plugins.}
  \begin{takeaway}{Green on your laptop proves nothing.}
    Run analyze and test on every pull request. A branch that fails the checks cannot be merged. The pipeline itself is Lecture 15.
  \end{takeaway}
  \note{Say that the CI job for this course runs the first two commands on every pull request, and that the lab grade depends on them passing. Integration tests need a device, so they usually run on a schedule rather than on every push. The integration\_test package comes from the SDK: add it under dev\_dependencies as integration\_test with sdk: flutter, never from pub.dev.}
\end{frame}

\begin{frame}[eyebrow=Recap]{Three things to remember}
  \vskip 2mm
  \begin{enumerate}
    \item Declare colors and type once. \muted{A widget that reads a role instead of a hex value already supports dark mode, large text and a rebrand.}
    \item Polish is mostly restraint. \muted{Short animations, honest error messages, a 48 dp target, and a screen that survives being made twice as wide.}
    \item Tests are how a refactor stays cheap. \muted{Many unit and widget tests, a few golden and integration tests, all of them run by CI.}
  \end{enumerate}
  \vskip 4mm
  \kv{Further reading}{\link{https://docs.flutter.dev/ui/design/material}{docs.flutter.dev/ui/design/material} · \link{https://docs.flutter.dev/ui/accessibility-and-internationalization}{Accessibility and i18n} · \link{https://docs.flutter.dev/testing/overview}{Testing overview}}
  \note{Next week is the last lecture: build modes, signing, the stores, push notifications, and the course recap. Remind them that the lab test rules are on that deck.}
\end{frame}

\closingframe{Questions?}{dinu\_stefan.rusu@upb.ro · Microsoft Teams: course channel}

\end{document}
